% Apunts de teoria en PDF · 3.º ESO · Unidad 3: Proporcionalidad y porcentajes
% Mateix estil que els apunts de Batxillerat (franja i caixes de teoria en blau, exemples en roig), però més visuals,
% curts i amb dibuixos, per a l'alumnat de 3r ESO. Es compila amb pdflatex dues vegades (per als enllaços de l'índex),
% amb logo-fentmates.pdf (_eines/logo/pdf/horitzontal-clar.pdf) al costat.
\documentclass[11pt,a4paper]{article}
\usepackage[utf8]{inputenc}
\usepackage[T1]{fontenc}
\usepackage[spanish,es-noshorthands,es-nodecimaldot,es-nolists]{babel}
\usepackage[sfdefault,book,lining]{FiraSans}
\usepackage[italic,noplusnominus,noequal]{mathastext}
\linespread{1.18}
\usepackage[a4paper,top=2.4cm,bottom=2.2cm,left=1.6cm,right=1.6cm,headheight=1.1cm,headsep=0.5cm,footskip=1cm]{geometry}
\usepackage{amsmath,amssymb,cancel}
\usepackage{icomma}
\usepackage{graphicx}
\usepackage[table]{xcolor}
\usepackage{tikz}\usetikzlibrary{calc,positioning,arrows.meta,decorations.pathreplacing,shapes.geometric}
\usepackage{fontawesome5}
\usepackage{ccicons}
\usepackage{fancyhdr}
\usepackage{enumitem}
\usepackage{array}
\usepackage{needspace}
\usepackage[most]{tcolorbox}
\usepackage{titlesec}
\usepackage[hidelinks]{hyperref}
\hypersetup{pdftitle={Apuntes de teoría. Unidad 3: Proporcionalidad y porcentajes. 3.º ESO},pdfsubject={Matemáticas 3.º ESO},pdfkeywords={porcentajes, índice de variación, interés compuesto, repartos proporcionales, mezclas, móviles, proporcionalidad simple, proporcionalidad compuesta}}

% colors de la web (bloc d'Àlgebra: blau)
\definecolor{fmnavy}{HTML}{1B2748}
\definecolor{fmblue}{HTML}{3D6AA0}
\definecolor{fmbluel}{HTML}{4E81BD}
\definecolor{fmyellow}{HTML}{FFC94D}
\definecolor{fmgold}{HTML}{E9A919}
\definecolor{fmaccent}{HTML}{FF7A50}
\definecolor{fmdef}{HTML}{1D4F73}\definecolor{fmdefbg}{HTML}{EEF4F9}
\definecolor{fmex}{HTML}{A3203A}\definecolor{fmexbg}{HTML}{FBF0F2}
\definecolor{fmsf}{HTML}{C2410C}
\definecolor{fmok}{HTML}{15803D}
\definecolor{fmmuted}{HTML}{5D6685}
\definecolor{fmsurface}{HTML}{EEF1F8}
\definecolor{fmA}{HTML}{2F6FD6}\definecolor{fmB}{HTML}{EA6A10}\definecolor{fmR}{HTML}{E11D48}\definecolor{fmK}{HTML}{16A34A}\definecolor{fmV}{HTML}{7C3AED}

% lletres
\newcommand{\comic}{\fontfamily{ComicNeue-TLF}\selectfont}
\newcommand{\heavy}{\fontseries{eb}\selectfont}
\newcommand{\semi}{\fontseries{sb}\selectfont}
\DeclareUnicodeCharacter{2212}{-}

\setlength{\parindent}{0pt}
\setlength{\parskip}{7pt}
\setlength{\lineskip}{4pt}\setlength{\lineskiplimit}{4pt}
\setlength{\emergencystretch}{3em}
\setlist{itemsep=3pt,topsep=3pt,leftmargin=1.4em}
\setcounter{secnumdepth}{1}
\renewcommand{\thesection}{3.\arabic{section}}
\titleformat{\section}[hang]{\color{fmnavy}\fontsize{23}{27}\heavy\selectfont}{\tcbox[on line,colback=fmblue,colframe=fmblue,arc=3mm,boxsep=2pt,left=6pt,right=6pt,top=3pt,bottom=3pt]{\color{white}\thesection}}{0.45em}{}[\vspace{5pt}{\color{fmblue}\titlerule[2pt]}]
\titlespacing*{\section}{0pt}{4pt}{6pt}
\titleformat{\subsection}{\color{fmblue}\Large\semi}{}{0pt}{\faCaretRight\ }
\titlespacing*{\subsection}{0pt}{8pt}{2pt}

% frase de cada apartat, «a mà»
\newcommand{\frase}[1]{\par\vspace{-2pt}{\comic\Large\color{fmbluel}\faCommentDots[regular]\ \ #1\par}\vspace{2pt}}
\newcommand{\seccio}[3]{\Needspace{16\baselineskip}\section{#2}\hypertarget{#1}{}\frase{#3}}
% notes a mà dins dels exemples
\newcommand{\nota}[1]{{\comic\small\color{fmex!85!black}#1}}
\newcommand{\notab}[1]{{\comic\small\color{fmblue}#1}}
% bé / malament
\newcommand{\be}[1]{{\color{fmok}\faCheckCircle}\ \ensuremath{#1}}
\newcommand{\mal}[1]{{\color{fmR}\faTimesCircle}\ {\color{fmR}\ensuremath{#1}}}
% pas numerat
\newcommand{\pas}[1]{\tikz[baseline=(p.base)]\node[circle,fill=fmblue,text=white,inner sep=1.2pt,minimum size=1.25em,font=\small\heavy](p){#1};}
% etiqueta de color
\newcommand{\chip}[2]{\tcbox[on line,colback=#1!14,colframe=#1!60,arc=2mm,boxsep=1pt,left=5pt,right=5pt,top=1pt,bottom=1pt]{#2}}
\newcommand{\fmkey}[1]{\colorbox{fmyellow!55}{\strut #1}}
\newcommand{\cA}[1]{{\color{fmA}#1}}\newcommand{\cB}[1]{{\color{fmB}#1}}\newcommand{\cK}[1]{{\color{fmK}#1}}\newcommand{\cR}[1]{{\color{fmR}#1}}\newcommand{\cV}[1]{{\color{fmV}#1}}

% caixes
\tcbset{fmbase/.style={enhanced,arc=3mm,boxrule=1.2pt,left=4mm,right=4mm,top=2mm,bottom=2.5mm,before skip=8pt,after skip=8pt,fonttitle=\heavy\large,before upper=\raggedright}}
\newtcolorbox{idea}[1][Idea clave]{fmbase,colback=fmdefbg,colframe=fmdef,coltitle=white,title={\faKey\ \ #1}}
\newtcolorbox{exemple}[1][Ejemplo]{fmbase,colback=fmexbg,colframe=fmex,colbacktitle=fmexbg,coltitle=fmex,title={\faPencil*\ \ #1},titlerule=0.5pt,titlerule style={fmex!40,dashed}}
\newtcolorbox{compte}[1][¡Cuidado!]{fmbase,colback=fmsf!6,colframe=fmsf,coltitle=white,fonttitle=\comic\bfseries\Large,title={\faExclamationTriangle\ \ #1}}
\newtcolorbox{truc}[1][Truco]{fmbase,colback=fmyellow!20,colframe=fmgold,colbacktitle=fmyellow,coltitle=fmnavy,fonttitle=\comic\bfseries\Large,title={\faLightbulb[regular]\ \ #1}}

% dibuix de la targeta de la unitat, en blanc
\newcommand{\fmcard}{\begin{tikzpicture}[x=0.062cm,y=-0.062cm,white,line width=3pt,line cap=round]
  \begin{scope}[transparency group,opacity=0.15]
  \node[anchor=base west,inner sep=0pt,font=\fontsize{34}{34}\selectfont] at (14,44) {$25\,\%\quad \frac{a}{b}=\frac{c}{d}$};
  \foreach \i in {0,...,4} \foreach \j in {0,...,4} {\draw[line width=1.6pt] (40+\i*10,60+\j*10) rectangle ++(10,10);}
  \fill[opacity=0.6] (40,60) rectangle (90,70);
  \draw[line width=3.2pt] (140,105) -- (140,55); \draw[line width=3.2pt] (140,105) -- (230,105);
  \draw[line width=3pt] (140,105) -- (220,62);
  \end{scope}
\end{tikzpicture}}
\newcommand{\fmbrand}{\raisebox{-0.25\height}{\includegraphics[height=0.5cm]{logo-fentmates.pdf}}}

% regla de Ruffini: #1 = a, #2 = coeficients, #3 = fila del mig (sense el primer), #4 = fila de baix, #5 = nombre de coeficients
\newcommand{\ruffini}[5]{\begin{tikzpicture}[x=1.25cm,y=0.85cm,every node/.style={font=\large}]
  \foreach \v [count=\i] in {#2} \node (t\i) at (\i,0) {$\v$};
  \foreach \v [count=\i] in {#3} {\pgfmathtruncatemacro\j{\i+1}\node (m\j) at (\j,-1) {$\v$};}
  \foreach \v [count=\i] in {#4} \node (b\i) at (\i,-2.1) {$\v$};
  \node[text=fmB] at (0,-1) {$#1$};
  \draw[line width=1.1pt,fmnavy] (0.45,0.5)--(0.45,-1.5) (0.45,-1.5)--(#5+0.5,-1.5);
  \node[draw=fmR,line width=1.3pt,rounded corners=2mm,minimum size=0.75cm] at (#5,-2.1) {};
  \draw[-{Latex[length=2.2mm]},fmK,line width=1.1pt] (t1.south) -- (b1.north);
  \pgfmathtruncatemacro\last{#5-1}
  \foreach \i in {1,...,\last} {\pgfmathtruncatemacro\j{\i+1}
    \draw[-{Latex[length=2.2mm]},fmB,line width=1.1pt,shorten >=1pt] (b\i.north east) -- (m\j.south west);
    \draw[-{Latex[length=2.2mm]},fmA!70,line width=1pt,dashed,shorten >=2pt] ([xshift=8pt]t\j.south) -- ([xshift=8pt]b\j.north);}
\end{tikzpicture}}

% capçalera i peu
\pagestyle{fancy}\fancyhf{}
\fancyhead[L]{\small\fmbrand}
\fancyhead[R]{\small\color{fmmuted}3.º ESO · Unidad 3: Proporcionalidad y porcentajes}
\fancyfoot[L]{\scriptsize\color{fmmuted}\href{https://creativecommons.org/licenses/by-nc-sa/4.0/deed.es}{\ccbyncsa}\ \ FentMates · fentmates.org · CC BY-NC-SA 4.0}
\fancyfoot[R]{\small\color{fmmuted}\thepage}
\renewcommand{\headrulewidth}{0.4pt}\renewcommand{\headrule}{\color{fmmuted!40}\hrule height \headrulewidth}
\fancypagestyle{primera}{\fancyhf{}\renewcommand{\headrulewidth}{0pt}
  \fancyfoot[L]{\scriptsize\color{fmmuted}\href{https://creativecommons.org/licenses/by-nc-sa/4.0/deed.es}{\ccbyncsa}\ \ FentMates · fentmates.org · CC BY-NC-SA 4.0}
  \fancyfoot[R]{\small\color{fmmuted}\thepage}}

% targeta de l'índex
\newcommand{\itile}[4]{\begin{tcolorbox}[enhanced,colback=white,colframe=fmblue!35,arc=3mm,boxrule=1pt,boxsep=0pt,left=6pt,right=6pt,top=5pt,bottom=5pt,before skip=0pt,after skip=7pt]
  \hyperlink{#1}{\begin{minipage}[c]{0.9cm}\centering{\color{fmblue}\LARGE #3}\end{minipage}\hspace{6pt}%
  \begin{minipage}[c]{0.78\linewidth}{\color{fmblue}\heavy #2}\\[-2pt]{\color{fmnavy}\semi\large #4}\end{minipage}}\end{tcolorbox}}

% balança
\newcommand{\balanca}[3]{\begin{tikzpicture}[scale=#3,transform shape,every node/.style={anchor=center}]
  \fill[fmnavy] (-0.45,-1.9) rectangle (0.45,-1.75);
  \fill[fmnavy!80] (-0.12,-1.75) rectangle (0.12,0);
  \fill[fmgold] (0,0.05) circle (0.13);
  \draw[fmnavy,line width=2.4pt,line cap=round] (-2,0) -- (2,0);
  \foreach \s in {-1,1} {
    \draw[fmmuted,line width=0.8pt] (\s*2,0) -- (\s*2-0.75,-1.05) (\s*2,0) -- (\s*2+0.75,-1.05);
    \fill[fmyellow!70,draw=fmgold,line width=1pt] (\s*2-0.95,-1.05) arc[start angle=180,end angle=360,x radius=0.95,y radius=0.32] -- cycle;}
  \node[font=\Large,fill=fmA!15,draw=fmA,rounded corners=1.5mm,inner sep=3pt] at (-2,-0.68) {$#1$};
  \node[font=\Large,fill=fmB!15,draw=fmB,rounded corners=1.5mm,inner sep=3pt] at (2,-0.68) {$#2$};
\end{tikzpicture}}
% eixos amb quadrícula (dins d'un tikzpicture): xmin, xmax, ymin, ymax
\newcommand{\eixos}[4]{\draw[fmmuted!22,line width=0.4pt] (#1,#3) grid (#2,#4);
  \draw[-{Latex[length=2mm]},fmmuted,line width=0.8pt] (#1-0.2,0)--(#2+0.4,0) node[right,font=\small]{$x$};
  \draw[-{Latex[length=2mm]},fmmuted,line width=0.8pt] (0,#3-0.2)--(0,#4+0.4) node[above,font=\small]{$y$};}
% taula de traduccions (diccionari) amb les vores arrodonides
\newcommand{\dic}[2]{#1 & {\large$#2$}\\}
\newcommand{\dicratio}{0.62}
\newcommand{\dichead}{\begin{tabular}{@{}>{\raggedright}p{\dimexpr\dicratio\linewidth-12pt}>{\centering\arraybackslash}p{\dimexpr\linewidth-\dicratio\linewidth-12pt}@{}}\faComment[regular]\ Lo decimos así\dots & \dots\ así\end{tabular}}
\newenvironment{dicbox}{\begin{tcolorbox}[enhanced,arc=3mm,boxrule=1.2pt,colframe=fmblue,colback=white,coltitle=white,fonttitle=\heavy,
  title={\dichead},left=0pt,right=0pt,top=0pt,bottom=1.5mm,boxsep=0pt,lefttitle=6pt,righttitle=6pt,before skip=0pt,after skip=0pt]\rowcolors{1}{white}{fmsurface}%
  \begin{tabular}{@{\hspace{6pt}}>{\raggedright}p{\dimexpr\dicratio\linewidth-12pt}>{\centering\arraybackslash}p{\dimexpr\linewidth-\dicratio\linewidth-12pt}@{\hspace{6pt}}}}
  {\end{tabular}\end{tcolorbox}}
% resum
\newcommand{\rtile}[3]{\begin{tcolorbox}[enhanced,colback=white,colframe=fmblue!45,arc=3mm,boxrule=1pt,left=3mm,right=3mm,top=2.5mm,bottom=2.5mm,fontupper=\large,
  title={\tcbox[on line,colback=fmblue,colframe=fmblue,arc=2mm,boxsep=0pt,left=4pt,right=4pt,top=1pt,bottom=1pt]{\color{white}\small\heavy #1}\ \ #2},coltitle=fmnavy,colbacktitle=fmdefbg,fonttitle=\semi,toptitle=0.5mm,bottomtitle=0.5mm]#3\end{tcolorbox}}
% portada: #1 subtítol a mà
\newcommand{\portada}[1]{\thispagestyle{primera}
\begin{tikzpicture}[remember picture,overlay]
  \fill[fmblue] (current page.north west) rectangle ([yshift=-9.4cm]current page.north east);
  \begin{scope}
    \clip (current page.north west) rectangle ([yshift=-9.4cm]current page.north east);
    \node[anchor=north west,inner sep=0pt] at ([xshift=6.6cm,yshift=-1.3cm]current page.north west) {\fmcard};
  \end{scope}
  \node[anchor=north west,inner sep=0pt] at ([xshift=1.6cm,yshift=-1.1cm]current page.north west) {\tcbox[colback=white,colframe=white,arc=2.5mm,boxsep=0pt,left=7pt,right=7pt,top=5pt,bottom=5pt]{\includegraphics[height=0.75cm]{logo-fentmates.pdf}}};
\end{tikzpicture}
\vspace*{0.55cm}
\tcbox[colback=white!15!fmblue,colframe=white!15!fmblue,arc=2.5mm,boxsep=0pt,left=6pt,right=6pt,top=3pt,bottom=3pt]{\color{white}\small\bfseries Apuntes de teoría · 3.º ESO · Matemáticas}

\vspace{2pt}
{\raggedright\color{white}\fontsize{30}{36}\heavy\selectfont{\color{fmyellow}Unidad 3}\\[2pt]Proporcionalidad y porcentajes\par}
\vspace{8pt}
{\color{white}\comic\Large #1\par}
\vspace*{\dimexpr 7.4cm-\pagetotal\relax}
}
\newcommand{\llegenda}{\par\vspace{16pt}
{\color{fmnavy}\heavy\Large\faPalette\ \ Cómo se leen estos apuntes}\par\smallskip
\begin{tcbraster}[raster columns=4,raster equal height,raster column skip=8pt,raster before skip=12pt]
\tcbox[colback=fmdefbg,colframe=fmdef,arc=3mm,boxrule=1.2pt,halign=center,width=\linewidth]{{\color{fmdef}\Large\faKey}\\[2pt]{\color{fmdef}\heavy Idea clave}\\[-1pt]{\small lo esencial}}
\tcbox[colback=fmexbg,colframe=fmex,arc=3mm,boxrule=1.2pt,halign=center,width=\linewidth]{{\color{fmex}\Large\faPencil*}\\[2pt]{\color{fmex}\heavy Ejemplo}\\[-1pt]{\small paso a paso}}
\tcbox[colback=fmsf!6,colframe=fmsf,arc=3mm,boxrule=1.2pt,halign=center,width=\linewidth]{{\color{fmsf}\Large\faExclamationTriangle}\\[2pt]{\color{fmsf}\heavy ¡Cuidado!}\\[-1pt]{\small errores típicos}}
\tcbox[colback=fmyellow!20,colframe=fmgold,arc=3mm,boxrule=1.2pt,halign=center,width=\linewidth]{{\color{fmgold!80!black}\Large\faLightbulb[regular]}\\[2pt]{\color{fmnavy}\heavy Truco}\\[-1pt]{\small para recordarlo}}
\end{tcbraster}}
\newcommand{\resumtitol}{\clearpage
\hypertarget{resum}{}
{\color{fmnavy}\fontsize{23}{27}\heavy\selectfont\faClipboardList\ \ En resumen}\par\vspace{5pt}{\color{fmblue}\titlerule[2pt]}
\frase{Toda la unidad en una página.}}
% els 5 passos per a resoldre problemes: #2 = text del pas 2
\newcommand{\passosproblema}[2]{\begin{center}
\begin{tikzpicture}[s/.style={draw=fmblue,fill=fmdefbg,rounded corners=3mm,line width=1.2pt,minimum width=2.75cm,minimum height=2.45cm,align=center,font=\small},
  ic/.style={font=\LARGE,text=fmblue}]
  \foreach \i/\ic/\t in {1/\faBookOpen/{\textbf{Leemos}\\y entendemos},2/\faQuestionCircle/{#1},3/\faEquals/{#2},4/\faCogs/{\textbf{Resolvemos}},5/\faCheckDouble/{\textbf{Comprobamos}\\y respondemos}} {
    \node[s] (s\i) at (3.3*\i,0) {};
    \node[ic] at ([yshift=-0.5cm]s\i.north) {\ic};
    \node[align=center,font=\small,anchor=north] at ([yshift=-0.95cm]s\i.north) {\t};
    \node[circle,fill=fmblue,text=white,font=\small\heavy,inner sep=1.5pt,minimum size=1.3em] at (s\i.north west) {\i};}
  \foreach \i [evaluate=\i as \j using int(\i+1)] in {1,...,4} \draw[-{Latex[length=3mm]},fmmuted,line width=1.5pt] (s\i) -- (s\j);
\end{tikzpicture}
\end{center}}

\newcommand{\pc}{\,\%}
% quadrícula de 100 amb #1 caselles pintades
\newcommand{\cent}[2]{\begin{tikzpicture}[x=#2,y=#2]
  \foreach \k in {0,...,99} {\pgfmathtruncatemacro\i{mod(\k,10)}\pgfmathtruncatemacro\j{9-int(\k/10)}
    \ifnum\k<#1 \filldraw[fill=fmB!45,draw=white,line width=0.6pt] (\i,\j) rectangle ++(1,1);
    \else \filldraw[fill=fmsurface,draw=white,line width=0.6pt] (\i,\j) rectangle ++(1,1);\fi}
  \draw[fmB,line width=1pt,rounded corners=1pt] (0,0) rectangle (10,10);
\end{tikzpicture}}
% capçalera de taula arrodonida (taules senzilles)
% taula amb vores arrodonides: el contingut es retalla amb la forma
\newsavebox\tabclipbox
\newcommand{\taulaclip}[3][3mm]{\sbox\tabclipbox{#3}\begin{tikzpicture}[baseline=(tcn.base)]
  \begin{scope}\clip[rounded corners=#1] (0,-\dp\tabclipbox) rectangle (\wd\tabclipbox,\ht\tabclipbox);
  \node[inner sep=0pt,outer sep=0pt,anchor=base west] (tcn) at (0,0) {\usebox\tabclipbox};\end{scope}
  \draw[#2,rounded corners=#1] (0,-\dp\tabclipbox) rectangle (\wd\tabclipbox,\ht\tabclipbox);
\end{tikzpicture}}
\newcommand{\taularod}[3][fmblue]{\taulaclip{#1,line width=1.2pt}{\renewcommand{\arraystretch}{1.5}\begin{tabular}{#2}#3\end{tabular}}}

\begin{document}
\portada{Porcentajes, repartos, mezclas\dots: las matemáticas de cada día.}

{\color{fmnavy}\heavy\Large\faMapSigns\ \ ¿Qué aprenderás?}\par\smallskip
\begin{minipage}[t]{0.485\linewidth}
\itile{sec1}{3.1}{\faPercent}{Porcentajes y fracciones}
\itile{sec2}{3.2}{\faCalculator}{Cálculo de porcentajes}
\itile{sec3}{3.3}{\faQuestionCircle}{¿Qué tanto por ciento es?}
\itile{sec4}{3.4}{\faArrowsAltV}{Aumentos y disminuciones}
\itile{sec5}{3.5}{\faUndo}{La cantidad inicial}
\itile{sec6}{3.6}{\faLink}{Variaciones encadenadas}
\itile{sec7}{3.7}{\faPiggyBank}{Interés compuesto}
\end{minipage}\hfill
\begin{minipage}[t]{0.485\linewidth}
\itile{sec8}{3.8}{\faUsers}{Repartos proporcionales}
\itile{sec9}{3.9}{\faBlender}{Problemas de mezclas}
\itile{sec10}{3.10}{\faCarSide}{Problemas de móviles}
\itile{sec11}{3.11}{\faBalanceScale}{Proporcionalidad simple}
\itile{sec12}{3.12}{\faCubes}{Proporcionalidad compuesta}
\itile{resum}{\phantom{3.12}}{\faClipboardList}{En resumen}
\end{minipage}
\llegenda

% ============================================================ 3.1
\clearpage
\seccio{sec1}{Porcentajes, fracciones y decimales}{Un porcentaje dice cuántas partes tomamos de cada 100.}

\begin{minipage}[c]{0.3\linewidth}\centering
\cent{35}{0.42cm}\\[4pt]
{\comic\color{fmB}35 de cada 100}
\end{minipage}\hfill
\begin{minipage}[c]{0.67\linewidth}
\begin{idea}[Tres maneras de decir lo mismo]
{\centering\fontsize{20}{26}\selectfont$\cB{35\pc}=\cA{\dfrac{35}{100}}=\cA{\dfrac{7}{20}}=\cK{0{,}35}$\par}\smallskip
{\centering\chip{fmB}{porcentaje}\quad \chip{fmA}{fracción}\quad \chip{fmK}{decimal (tanto por uno)}\par}\smallskip
El $100\pc$ es el \textbf{total} ($=1$). Más del $100\pc$ es \textbf{más} que el total.
\end{idea}
\end{minipage}

\subsection{Cómo pasamos de una forma a otra}
\begin{center}
\begin{tikzpicture}[b/.style={draw=#1,fill=#1!12,line width=1.3pt,rounded corners=3mm,minimum width=3.6cm,minimum height=1.3cm,align=center,font=\large}]
  \node[b=fmB] (p) at (0,0) {\heavy Porcentaje\\$35\pc$};
  \node[b=fmK] (d) at (7,0) {\heavy Decimal\\$0{,}35$};
  \node[b=fmA] (f) at (3.5,-2.6) {\heavy Fracción\\$\dfrac{35}{100}=\dfrac{7}{20}$};
  \draw[-{Latex[length=3mm]},line width=1.3pt,fmmuted] ([yshift=5pt]p.east) -- node[above,font=\comic]{entre 100} ([yshift=5pt]d.west);
  \draw[-{Latex[length=3mm]},line width=1.3pt,fmmuted] ([yshift=-5pt]d.west) -- node[below,font=\comic]{por 100} ([yshift=-5pt]p.east);
  \draw[-{Latex[length=3mm]},line width=1.3pt,fmmuted] (p.south) -- node[left=4pt,font=\comic,align=right]{$\dfrac{p}{100}$ y\\simplificamos} (f.west);
  \draw[-{Latex[length=3mm]},line width=1.3pt,fmmuted] (f.east) -- node[right=4pt,font=\comic,align=left]{dividimos\\arriba entre abajo} (d.south);
\end{tikzpicture}
\end{center}

\begin{exemple}[Ejemplos]
\begin{minipage}[t]{0.48\linewidth}
$4\pc=0{,}04\qquad 120\pc=1{,}2$\\[3pt]
$0{,}375=37{,}5\pc\qquad 1{,}35=135\pc$
\end{minipage}\hfill
\begin{minipage}[t]{0.49\linewidth}
$60\pc=\dfrac{60}{100}=\dfrac35$\\[3pt]
$\dfrac38=0{,}375=37{,}5\pc\qquad \dfrac{7}{25}=\dfrac{28}{100}=28\pc$
\end{minipage}
\end{exemple}

\begin{center}
\taularod{>{\columncolor{fmblue}\color{white}\heavy}l*{9}{w{c}{1.08cm}}}{
\rowcolor{fmblue}\% & \color{white}$1\pc$ & \color{white}$5\pc$ & \color{white}$10\pc$ & \color{white}$20\pc$ & \color{white}$25\pc$ & \color{white}$33{,}\overline{3}\pc$ & \color{white}$50\pc$ & \color{white}$75\pc$ & \color{white}$100\pc$\\
Fracción & $\frac{1}{100}$ & $\frac{1}{20}$ & $\frac{1}{10}$ & $\frac15$ & $\frac14$ & $\frac13$ & $\frac12$ & $\frac34$ & $1$\\
\rowcolor{fmsurface}\cellcolor{fmblue}Decimal & $0{,}01$ & $0{,}05$ & $0{,}1$ & $0{,}2$ & $0{,}25$ & $0{,}\overline{3}$ & $0{,}5$ & $0{,}75$ & $1$
}
\end{center}
{\comic\color{fmblue}\centering Vale la pena saberla de memoria: el $50\pc$ es la mitad, el $25\pc$ la cuarta parte, el $10\pc$ la décima parte\dots\par}

% ============================================================ 3.2
\clearpage
\seccio{sec2}{Cálculo de porcentajes}{Multiplicamos la cantidad por el porcentaje en forma decimal.}

\begin{minipage}[c]{0.52\linewidth}
\begin{tcolorbox}[enhanced,colback=fmyellow!18,colframe=fmgold,arc=4mm,boxrule=1.6pt,before upper=\centering,top=3mm,bottom=3mm]
{\Large$p\pc$ de $C=C\cdot\dfrac{p}{100}$}\par\medskip
$15\pc$ de $80=80\cdot0{,}15=\boxed{12}$
\end{tcolorbox}
\end{minipage}\hfill
\begin{minipage}[c]{0.45\linewidth}
\begin{idea}[Es una regla de tres]
{\centering$\left.\begin{array}{l}100\ \longrightarrow\ p\\ C\ \ \ \longrightarrow\ x\end{array}\right\}\quad x=\dfrac{C\cdot p}{100}$\par}
\end{idea}
\end{minipage}

\subsection{Trucos de cálculo mental}
\begin{tcbraster}[raster columns=4,raster equal height,raster column skip=8pt]
\begin{tcolorbox}[fmbase,colback=fmdefbg,colframe=fmdef,before upper=\centering,top=2mm,bottom=2mm]
{\Large\heavy\color{fmblue}$50\pc$}\par la \textbf{mitad}\par{\small$50\pc$ de $84=42$}
\end{tcolorbox}
\begin{tcolorbox}[fmbase,colback=fmdefbg,colframe=fmdef,before upper=\centering,top=2mm,bottom=2mm]
{\Large\heavy\color{fmblue}$25\pc$}\par la mitad de la mitad\par{\small$25\pc$ de $84=21$}
\end{tcolorbox}
\begin{tcolorbox}[fmbase,colback=fmdefbg,colframe=fmdef,before upper=\centering,top=2mm,bottom=2mm]
{\Large\heavy\color{fmblue}$10\pc$}\par entre 10\par{\small$10\pc$ de $84=8{,}4$}
\end{tcolorbox}
\begin{tcolorbox}[fmbase,colback=fmdefbg,colframe=fmdef,before upper=\centering,top=2mm,bottom=2mm]
{\Large\heavy\color{fmblue}$1\pc$}\par entre 100\par{\small$1\pc$ de $84=0{,}84$}
\end{tcolorbox}
\end{tcbraster}

\begin{minipage}[t]{0.5\linewidth}
\begin{truc}[Combínalos]
$5\pc$: la mitad del $10\pc$. \ $20\pc$: el doble del $10\pc$.\\[4pt]
$15\pc$ de $60=\underbrace{6}_{10\pc}+\underbrace{3}_{5\pc}=9$
\end{truc}
\end{minipage}\hfill
\begin{minipage}[t]{0.47\linewidth}
\begin{truc}[Se pueden intercambiar]
El $p\pc$ de $C$ es igual que el $C\pc$ de $p$:\\[4pt]
el $8\pc$ de $50$ = el $50\pc$ de $8=\boxed{4}$
\end{truc}
\end{minipage}

\begin{exemple}[Ejemplo: \ de 640 alumnos, el $45\pc$ son chicas. ¿Cuántas hay?]
$45\pc$ de $640=640\cdot0{,}45=\boxed{288}$ chicas \qquad \nota{o: $\ 10\pc=64,\ \ 40\pc=256,\ \ 5\pc=32\ \to\ 256+32=288$}
\end{exemple}

% ============================================================ 3.3
\seccio{sec3}{¿Qué tanto por ciento representa?}{Dividimos la parte entre el total y multiplicamos por 100.}

\begin{minipage}[c]{0.45\linewidth}
\begin{tcolorbox}[enhanced,colback=fmyellow!18,colframe=fmgold,arc=4mm,boxrule=1.6pt,before upper=\centering,top=3mm,bottom=3mm]
{\Large$p=\dfrac{\cB{P}}{\cA{C}}\cdot100$}\par\smallskip
{\small\cB{parte} entre \cA{total}}\par\smallskip
18 de 24: \ $\dfrac{18}{24}=0{,}75=\boxed{75\pc}$
\end{tcolorbox}
\end{minipage}\hfill
\begin{minipage}[c]{0.52\linewidth}
\taularod{>{\heavy}l<{\hspace{5pt}}>{\hspace{5pt}}l<{\hspace{5pt}}>{\hspace{5pt}}l<{\hspace{4pt}}}{
\rowcolor{fmblue}\color{white}Queremos & \color{white}\heavy Cálculo & \color{white}\heavy Ejemplo\\
la parte $P$ & $C\cdot\frac{p}{100}$ & $20\pc$ de $45=9$\\
\rowcolor{fmsurface}el $\%$ & $\frac{P}{C}\cdot100$ & 9 de 45 $=20\pc$\\
el total $C$ & $P:\frac{p}{100}$ & 9 es el $20\pc$ de $45$
}
\end{minipage}

% ============================================================ 3.4
\clearpage
\seccio{sec4}{Aumentos y disminuciones porcentuales}{Multiplicamos directamente por el índice de variación.}

\begin{center}
\begin{tikzpicture}[x=0.075cm]
  \filldraw[fill=fmA!25,draw=fmA,line width=1.2pt] (0,0) rectangle (100,0.8); \node at (50,0.4) {\large$100\pc$ \ (inicial)};
  \filldraw[fill=fmK!30,draw=fmK,line width=1.2pt] (100,0) rectangle (115,0.8); \node at (107.5,0.4) {$+15$};
  \draw[decorate,decoration={brace,amplitude=5pt},fmK,line width=1.1pt] (0,0.95) -- (115,0.95) node[midway,above=6pt,font=\large,text=fmK!80!black]{aumento del $15\pc$: \ $115\pc\ \to\ I_V=1{,}15$};
  \begin{scope}[yshift=-1.9cm]
  \filldraw[fill=fmA!25,draw=fmA,line width=1.2pt] (0,0) rectangle (80,0.8); \node at (40,0.4) {\large$80\pc$ \ (queda)};
  \filldraw[fill=fmR!8,draw=fmR,line width=1.2pt,dashed] (80,0) rectangle (100,0.8); \node[text=fmR] at (90,0.4) {$-20$};
  \draw[decorate,decoration={brace,amplitude=5pt,mirror},fmR,line width=1.1pt] (0,-0.15) -- (80,-0.15) node[midway,below=6pt,font=\large,text=fmR]{disminución del $20\pc$: \ $80\pc\ \to\ I_V=0{,}8$};
  \end{scope}
\end{tikzpicture}
\end{center}

\begin{tcolorbox}[enhanced,colback=fmyellow!18,colframe=fmgold,arc=4mm,boxrule=1.6pt,before upper=\centering,top=3mm,bottom=3mm]
{\Large$C_f=C_i\cdot I_V$}\qquad\qquad aumento: \ $I_V=1+\dfrac{p}{100}$\qquad\quad disminución: \ $I_V=1-\dfrac{p}{100}$
\end{tcolorbox}

\begin{minipage}[t]{0.5\linewidth}
\begin{exemple}[Ejemplos]
Un precio de 80 € sube un $15\pc$:\\[2pt]
\hspace*{1em}$80\cdot1{,}15=\boxed{92\text{ €}}$\\[4pt]
Una chaqueta de 50 € tiene un $20\pc$ de descuento:\\[2pt]
\hspace*{1em}$50\cdot0{,}8=\boxed{40\text{ €}}$
\end{exemple}
\end{minipage}\hfill
\begin{minipage}[t]{0.47\linewidth}
\begin{idea}[Leer un índice de variación]
$I_V>1$: \textbf{aumenta}. \ $I_V<1$: \textbf{disminuye}.\\[4pt]
$1{,}15\to\cK{+15\pc}\qquad 0{,}92\to\cR{-8\pc}$\\[3pt]
$1{,}05\to\cK{+5\pc}\qquad\ \ 0{,}7\to\cR{-30\pc}$\\[3pt]
$2\to\cK{+100\pc}$
\end{idea}
\end{minipage}

\begin{minipage}[t]{0.5\linewidth}
\begin{truc}[¿Conoces el inicio y el final?]
$I_V=\dfrac{C_f}{C_i}$: \ de 40 € a 46 €, \ $\dfrac{46}{40}=1{,}15\to+15\pc$
\end{truc}
\end{minipage}\hfill
\begin{minipage}[t]{0.47\linewidth}
\begin{compte}[Un solo paso]
No hace falta calcular el aumento y sumarlo: multiplica por $I_V$ y ya lo tienes.
\end{compte}
\end{minipage}

\begin{exemple}[Más ejemplos]
\begin{minipage}[t]{0.48\linewidth}
\textbf{IVA del $21\pc$}: un artículo vale 50 € sin IVA.\\[3pt]
$50\cdot1{,}21=\boxed{60{,}50\text{ €}}$
\end{minipage}\hfill
\begin{minipage}[t]{0.49\linewidth}
Un pueblo de 25\,000 habitantes pierde el $4\pc$.\\[3pt]
$25\,000\cdot0{,}96=\boxed{24\,000\text{ habitantes}}$
\end{minipage}
\end{exemple}

% ============================================================ 3.5
\clearpage
\seccio{sec5}{Cálculo de la cantidad inicial}{Para volver atrás, dividimos entre el índice de variación.}

\begin{center}
\begin{tikzpicture}
  \node[draw=fmA,fill=fmA!12,line width=1.3pt,rounded corners=3mm,minimum width=3cm,minimum height=1.2cm,font=\Large] (i) at (0,0) {$C_i$};
  \node[draw=fmK,fill=fmK!12,line width=1.3pt,rounded corners=3mm,minimum width=3cm,minimum height=1.2cm,font=\Large] (f) at (8,0) {$C_f$};
  \draw[-{Latex[length=3mm]},line width=1.6pt,fmK] ([yshift=6pt]i.east) to[bend left=18] node[above,font=\Large]{$\cdot\,I_V$} ([yshift=6pt]f.west);
  \draw[-{Latex[length=3mm]},line width=1.6pt,fmB] ([yshift=-6pt]f.west) to[bend left=18] node[below,font=\Large]{$:\,I_V$} ([yshift=-6pt]i.east);
\end{tikzpicture}
\end{center}

\begin{minipage}[c]{0.4\linewidth}
\begin{tcolorbox}[enhanced,colback=fmyellow!18,colframe=fmgold,arc=4mm,boxrule=1.6pt,before upper=\centering,top=3mm,bottom=3mm]
{\Large$C_i=\dfrac{C_f}{I_V}$}
\end{tcolorbox}
\end{minipage}\hfill
\begin{minipage}[c]{0.57\linewidth}
\begin{compte}[Error muy frecuente]
Aplicar el porcentaje a la cantidad \textbf{final}. ¡El \% se calculaba sobre la \textbf{inicial}!
\end{compte}
\end{minipage}

\begin{minipage}[t]{0.49\linewidth}
\begin{exemple}[Ejemplo: \ con un $-25\pc$ pagamos 60 €]
$I_V=0{,}75$ \ $\to$ \ $60:0{,}75=\boxed{80\text{ €}}$\\[3pt]
\mal{60+25\pc\text{ de }60=75\text{ €}}\\[2pt]
\nota{comprobación: $80\cdot0{,}75=60$ \be{}}
\end{exemple}
\end{minipage}\hfill
\begin{minipage}[t]{0.48\linewidth}
\begin{exemple}[Ejemplo: \ tras subir un $12\pc$ vale 56 €]
$I_V=1{,}12$ \ $\to$ \ $56:1{,}12=\boxed{50\text{ €}}$\\[3pt]
\nota{comprobación: $50\cdot1{,}12=56$ \be{}}
\end{exemple}
\end{minipage}

% ============================================================ 3.6
\seccio{sec6}{Variaciones porcentuales encadenadas}{Cada variación se aplica sobre el resultado de la anterior.}

\begin{minipage}[c]{0.48\linewidth}
\begin{tcolorbox}[enhanced,colback=fmyellow!18,colframe=fmgold,arc=4mm,boxrule=1.6pt,before upper=\centering,top=3mm,bottom=3mm]
{\Large$I_{V\,\text{total}}=I_{V_1}\cdot I_{V_2}\cdot I_{V_3}\cdots$}\par\smallskip
los índices se \textbf{multiplican}
\end{tcolorbox}
\end{minipage}\hfill
\begin{minipage}[c]{0.49\linewidth}
\begin{compte}[¡NO se suman!]
Sube un $20\pc$ y después baja un $20\pc$:\\[3pt]
$1{,}2\cdot0{,}8=0{,}96\ \to\ \cR{-4\pc}$\\[2pt]
\notab{¡No volvemos al precio inicial!}
\end{compte}
\end{minipage}

\begin{exemple}[Ejemplo: \ un precio de 200 € sube un $10\pc$, después baja un $30\pc$ y después sube un $5\pc$]
$I_V=1{,}1\cdot0{,}7\cdot1{,}05=0{,}8085$ \quad $\to$ \quad $200\cdot0{,}8085=\boxed{161{,}70\text{ €}}$ \quad \nota{en total, baja un $19{,}15\pc$}
\end{exemple}

% ============================================================ 3.7
\clearpage
\seccio{sec7}{Interés compuesto}{Los intereses también producen intereses.}

\begin{tcbraster}[raster columns=2,raster equal height,raster column skip=10pt]
\begin{tcolorbox}[fmbase,colback=fmdefbg,colframe=fmdef,coltitle=white,title={Interés simple}]
Los intereses se calculan siempre sobre el capital \textbf{inicial} y se retiran.\par\smallskip
{\centering$I=\dfrac{C_0\cdot r\cdot t}{100}$\par}
\end{tcolorbox}
\begin{tcolorbox}[fmbase,colback=fmK!8,colframe=fmK,coltitle=white,title={Interés compuesto}]
Cada año los intereses se \textbf{suman} al capital: es una variación encadenada.\par\smallskip
{\centering\Large$C_f=C_0\cdot\left(1+\dfrac{r}{100}\right)^{t}$\par}
\end{tcolorbox}
\end{tcbraster}
{\comic\color{fmmuted}\centering $C_0$: capital inicial \quad $r$: rédito (\% anual) \quad $t$: años \quad $I=C_f-C_0$: intereses\par}

\begin{exemple}[Ejemplo: \ depositamos 1000 € al $5\pc$ anual durante 3 años]
\begin{minipage}[c]{0.48\linewidth}\centering
\begin{tikzpicture}[x=1.15cm,y=0.0042cm]
  \draw[fmmuted,line width=0.8pt] (-0.4,0) -- (3.6,0);
  \foreach \t/\v/\l in {0/1000/1000,1/1050/1050,2/1102.5/{1102{,}50},3/1157.625/{1157{,}63}} {
    \filldraw[fill=fmK!35,draw=fmK,line width=1pt] (\t-0.32,0) rectangle (\t+0.32,\v-800);
    \node[above,font=\small] at (\t,\v-800) {\l};
    \node[below,font=\small] at (\t,0) {año \t};}
\end{tikzpicture}
\end{minipage}\hfill
\begin{minipage}[c]{0.49\linewidth}
\textbf{Compuesto}: \ $1000\cdot1{,}05^{3}=\boxed{1157{,}63\text{ €}}$\\[2pt]
\hspace*{1em}intereses: $157{,}63$ €\\[6pt]
\textbf{Simple}: \ $I=\dfrac{1000\cdot5\cdot3}{100}=150$ €\\[4pt]
\nota{Con el compuesto ganamos más: los intereses crecen cada año.}
\end{minipage}
\end{exemple}

\begin{minipage}[t]{0.55\linewidth}
\begin{idea}[Si se abona $n$ veces al año]
Mensual $n=12$, trimestral $n=4$, semestral $n=2$:\\[4pt]
{\centering\large$C_f=C_0\cdot\left(1+\dfrac{r}{100\cdot n}\right)^{n\cdot t}$\par}
\end{idea}
\end{minipage}\hfill
\begin{minipage}[t]{0.42\linewidth}
\begin{exemple}[Ejemplo]
1000 € al $6\pc$, trimestral, 2 años:\\[3pt]
$1000\cdot\left(1+\dfrac{6}{400}\right)^{8}=$\\[2pt]
$=1000\cdot1{,}015^{8}\approx\boxed{1126{,}49\text{ €}}$
\end{exemple}
\end{minipage}

\begin{exemple}[Ejemplo: \ 2000 € al $3\pc$ anual durante 4 años. ¿Cuántos intereses ganamos?]
\begin{minipage}[t]{0.5\linewidth}
\textbf{Compuesto}: \ $2000\cdot1{,}03^{4}\approx2251{,}02$ €\\[2pt]
\hspace*{1em}$I=2251{,}02-2000=\boxed{251{,}02\text{ €}}$
\end{minipage}\hfill
\begin{minipage}[t]{0.47\linewidth}
\textbf{Simple}: \ $I=\dfrac{2000\cdot3\cdot4}{100}=240$ €\\[2pt]
\nota{la diferencia crece cuantos más años pasan}
\end{minipage}
\end{exemple}

% ============================================================ 3.8
\clearpage
\seccio{sec8}{Repartos proporcionales}{Repartir de manera justa: según lo que aporta cada uno.}

\begin{tcbraster}[raster columns=2,raster equal height,raster column skip=10pt]
\begin{tcolorbox}[fmbase,colback=fmK!8,colframe=fmK,coltitle=white,title={\faArrowUp\ \ Directamente proporcional}]
\textbf{Más} número, \textbf{más} parte (el doble, el doble).\par\smallskip
\pas{1} Una parte: \ $k=\dfrac{N}{a+b+c}$\\[3pt]
\pas{2} Cada uno: \ $k\cdot a$, \ $k\cdot b$, \ $k\cdot c$\\[3pt]
\pas{3} Comprobamos que suman $N$.
\end{tcolorbox}
\begin{tcolorbox}[fmbase,colback=fmB!8,colframe=fmB,coltitle=white,title={\faArrowDown\ \ Inversamente proporcional}]
\textbf{Más} número, \textbf{menos} parte (el doble, la mitad).\par\smallskip
Es repartir directamente a los \textbf{inversos}: $\dfrac1a,\ \dfrac1b,\ \dfrac1c$\\[3pt]
{\centering$k=\dfrac{N}{\frac1a+\frac1b+\frac1c}\qquad x=\dfrac{k}{a}\dots$\par}
\end{tcolorbox}
\end{tcbraster}

\begin{exemple}[Ejemplo: \ repartimos 900 € directamente proporcional a 2, 3 y 4 (horas trabajadas)]
\begin{minipage}[c]{0.5\linewidth}\centering
\begin{tikzpicture}[x=0.7cm]
  \foreach \i/\c in {0/fmA,1/fmA,2/fmB,3/fmB,4/fmB,5/fmK,6/fmK,7/fmK,8/fmK} {
    \filldraw[fill=\c!35,draw=\c,line width=1pt] (\i,0) rectangle (\i+1,0.65);
    \node[font=\scriptsize] at (\i+0.5,0.32) {100};}
  \node[below,font=\small,text=fmA] at (1,0) {2 partes};
  \node[below,font=\small,text=fmB] at (3.5,0) {3 partes};
  \node[below,font=\small,text=fmK!80!black] at (7,0) {4 partes};
\end{tikzpicture}
\end{minipage}\hfill
\begin{minipage}[c]{0.47\linewidth}
$k=\dfrac{900}{2+3+4}=100$ € cada parte\\[4pt]
\cA{$200$ €}, \ \cB{$300$ €}, \ \cK{$400$ €}\\[2pt]
\nota{$200+300+400=900$ \be{}}
\end{minipage}
\end{exemple}

\begin{minipage}[t]{0.52\linewidth}
\begin{exemple}[Ejemplo: \ 600 € inversamente a 2, 3 y 6]
Inversos: \ $\dfrac12,\ \dfrac13,\ \dfrac16=\dfrac36,\ \dfrac26,\ \dfrac16$\\[4pt]
Directamente a \textbf{3, 2 y 1}: \ $k=\dfrac{600}{6}=100$\\[4pt]
$\boxed{300\text{ €}},\ \boxed{200\text{ €}},\ \boxed{100\text{ €}}$\\[2pt]
\nota{quien tiene el 2 (el menor) recibe más}
\end{exemple}
\end{minipage}\hfill
\begin{minipage}[t]{0.45\linewidth}
\begin{truc}[Sin fracciones]
Reducimos los inversos a \textbf{común denominador} (m.c.m.) y repartimos \textbf{directamente} a los numeradores.
\end{truc}
\end{minipage}

\begin{exemple}[Ejemplo: \ tres amigas ponen 10 €, 15 € y 25 € en un boleto y ganan 150 €]
\begin{minipage}[c]{0.5\linewidth}
Reparto \textbf{directo} a 10, 15 y 25:\\[3pt]
$k=\dfrac{150}{10+15+25}=\dfrac{150}{50}=3$
\end{minipage}\hfill
\begin{minipage}[c]{0.47\linewidth}
$10\cdot3=\boxed{30\text{ €}}$, \ $15\cdot3=\boxed{45\text{ €}}$, \ $25\cdot3=\boxed{75\text{ €}}$\\[2pt]
\nota{$30+45+75=150$ \be{}}
\end{minipage}
\end{exemple}

% ============================================================ 3.9
\clearpage
\seccio{sec9}{Problemas de mezclas}{El precio de la mezcla es el coste total entre la cantidad total.}

\begin{minipage}[c]{0.52\linewidth}
\begin{tcolorbox}[enhanced,colback=fmyellow!18,colframe=fmgold,arc=4mm,boxrule=1.6pt,before upper=\centering,top=3mm,bottom=3mm]
{\Large$p_m=\dfrac{m_1\cdot p_1+m_2\cdot p_2}{m_1+m_2}$}\par\smallskip
{\comic\color{fmmuted}coste total entre cantidad total}
\end{tcolorbox}
\end{minipage}\hfill
\begin{minipage}[c]{0.45\linewidth}
\begin{idea}
El precio de la mezcla queda \textbf{entre} los dos precios, y más cerca del que entra en \textbf{más cantidad}.
\end{idea}
\end{minipage}

\begin{exemple}[Ejemplo: \ mezclamos 30 kg de café de 8 €/kg con 20 kg de 13 €/kg]
\begin{minipage}[c]{0.55\linewidth}
$p_m=\dfrac{30\cdot8+20\cdot13}{30+20}=\dfrac{240+260}{50}=\boxed{10\text{ €/kg}}$
\end{minipage}\hfill
\begin{minipage}[c]{0.42\linewidth}
\nota{Queda entre 8 y 13, más cerca del 8 (hay más cantidad).}
\end{minipage}
\end{exemple}

\subsection{El método de la balanza}
\begin{minipage}[c]{0.45\linewidth}\centering
\begin{tikzpicture}[x=0.9cm]
  \draw[fmmuted,line width=1pt] (6,0) -- (12,0);
  \foreach \k in {6,8,12} {\draw[fmmuted] (\k,-0.1) -- (\k,0.1); \node[below,font=\small] at (\k,-0.1) {$\k$ €};}
  \fill[fmgold] (8,0.05) -- (7.75,-0.45) -- (8.25,-0.45) -- cycle;
  \draw[fmnavy,line width=2pt] (6,0.25) -- (12,0.25);
  \node[draw=fmA,fill=fmA!15,rounded corners=1.5mm,font=\small,anchor=south] at (6,0.35) {$x$ kg};
  \node[draw=fmB,fill=fmB!15,rounded corners=1.5mm,font=\small,anchor=south] at (12,0.35) {10 kg};
  \draw[decorate,decoration={brace,amplitude=4pt},line width=0.8pt,fmA] (6,1.05) -- (8,1.05) node[midway,above=4pt,font=\small,text=fmA]{$2$};
  \draw[decorate,decoration={brace,amplitude=4pt},line width=0.8pt,fmB] (8,1.05) -- (12,1.05) node[midway,above=4pt,font=\small,text=fmB]{$4$};
\end{tikzpicture}\\[2pt]
{\comic\small\color{fmmuted}cada cantidad «pesa» según la distancia al precio de la mezcla}
\end{minipage}\hfill
\begin{minipage}[c]{0.52\linewidth}
{\large$m_1\cdot(p_m-p_1)=m_2\cdot(p_2-p_m)$}\\[6pt]
\begin{exemple}[Ejemplo]
¿Cuántos kg de 6 €/kg mezclamos con 10 kg de 12 €/kg para obtener una mezcla de 8 €/kg?\\[3pt]
$x\cdot(8-6)=10\cdot(12-8)\ \to\ 2x=40\ \to\ \boxed{x=20\text{ kg}}$
\end{exemple}
\end{minipage}

\begin{minipage}[t]{0.5\linewidth}
\begin{idea}[Aleaciones: la ley]
{\centering$\text{ley}=\dfrac{\text{masa de metal precioso}}{\text{masa total}}$\par}\smallskip
Un número entre 0 y 1. También sirve para concentraciones (\% de zumo, de sal\dots).
\end{idea}
\end{minipage}\hfill
\begin{minipage}[t]{0.47\linewidth}
\begin{truc}[Los quilates]
El oro puro tiene \textbf{24 quilates}. Un oro de 18 k tiene ley \ $\dfrac{18}{24}=0{,}750$.\\[3pt]
\notab{Una pieza de 40 g de oro de 18 k tiene $40\cdot0{,}75=30$ g de oro puro.}
\end{truc}
\end{minipage}

\begin{exemple}[Ejemplo: \ mezclamos 2 L de zumo al $40\pc$ con 3 L al $20\pc$. ¿Qué concentración tiene la mezcla?]
$\dfrac{2\cdot0{,}40+3\cdot0{,}20}{2+3}=\dfrac{0{,}8+0{,}6}{5}=0{,}28\ \to\ \boxed{28\pc}$ \qquad \nota{entre el 20 y el 40, más cerca del 20}
\end{exemple}

% ============================================================ 3.10
\clearpage
\seccio{sec10}{Problemas de móviles}{Distancia = velocidad $\times$ tiempo.}

\begin{minipage}[c]{0.45\linewidth}
\begin{tcolorbox}[enhanced,colback=fmyellow!18,colframe=fmgold,arc=4mm,boxrule=1.6pt,before upper=\centering,top=3mm,bottom=3mm]
{\fontsize{24}{28}\selectfont$d=v\cdot t$}\par\medskip
$v=\dfrac{d}{t}\qquad t=\dfrac{d}{v}$
\end{tcolorbox}
\end{minipage}\hfill
\begin{minipage}[c]{0.52\linewidth}
\begin{compte}[¡Cuidado con las unidades!]
Con km/h, $t$ en \textbf{horas}: \ 2 h 15 min $=2{,}25$ h.\\[4pt]
De km/h a m/s, \textbf{dividimos entre 3,6}:\\[2pt]
$72$ km/h $=\dfrac{72\,000\text{ m}}{3\,600\text{ s}}=20$ m/s
\end{compte}
\end{minipage}

\subsection{Dos móviles}
\begin{tcbraster}[raster columns=2,raster equal height,raster column skip=10pt]
\begin{tcolorbox}[fmbase,colback=fmdefbg,colframe=fmdef,coltitle=white,title={\faArrowRight\,\faArrowLeft\ \ Se encuentran (de frente)}]
{\centering\begin{tikzpicture}[x=0.6cm]
  \draw[fmmuted,line width=1pt] (0,0) -- (10,0);
  \node[text=fmA,font=\Large] at (0.4,0.4) {\faCarSide}; \node[text=fmB,font=\Large,xscale=-1] at (9.6,0.4) {\faCarSide};
  \draw[-{Latex},fmA,line width=1.2pt] (1.3,0.4) -- (3,0.4); \draw[-{Latex},fmB,line width=1.2pt] (8.7,0.4) -- (7,0.4);
  \draw[decorate,decoration={brace,amplitude=4pt,mirror},line width=0.8pt] (0,-0.15) -- (10,-0.15) node[midway,below=4pt,font=\small]{$d$};
\end{tikzpicture}\par}\smallskip
La distancia se reduce a $v_1+v_2$:\par\smallskip
{\centering\large$t=\dfrac{d}{v_1+v_2}$\par}
\end{tcolorbox}
\begin{tcolorbox}[fmbase,colback=fmdefbg,colframe=fmdef,coltitle=white,title={\faArrowRight\,\faArrowRight\ \ Lo alcanza (mismo sentido)}]
{\centering\begin{tikzpicture}[x=0.6cm]
  \draw[fmmuted,line width=1pt] (0,0) -- (10,0);
  \node[text=fmA,font=\Large] at (0.4,0.4) {\faCarSide}; \node[text=fmB,font=\Large] at (5.4,0.4) {\faCarSide};
  \draw[-{Latex},fmA,line width=1.2pt] (1.3,0.4) -- (3.4,0.4); \draw[-{Latex},fmB,line width=1.2pt] (6.3,0.4) -- (7.5,0.4);
  \draw[decorate,decoration={brace,amplitude=4pt,mirror},line width=0.8pt] (0,-0.15) -- (5,-0.15) node[midway,below=4pt,font=\small]{$d$};
\end{tikzpicture}\par}\smallskip
La distancia se reduce a $v_1-v_2$:\par\smallskip
{\centering\large$t=\dfrac{d}{v_1-v_2}$\par}
\end{tcolorbox}
\end{tcbraster}

\begin{minipage}[t]{0.49\linewidth}
\begin{exemple}[Ejemplo: \ se encuentran]
Dos ciudades están a 300 km. Salen a la vez uno hacia el otro a 80 km/h y a 70 km/h.\\[3pt]
$t=\dfrac{300}{80+70}=\dfrac{300}{150}=\boxed{2\text{ h}}$\\[2pt]
\nota{el primero ha hecho $80\cdot2=160$ km}
\end{exemple}
\end{minipage}\hfill
\begin{minipage}[t]{0.48\linewidth}
\begin{exemple}[Ejemplo: \ lo alcanza]
Un camión va a 60 km/h y lleva 30 km de ventaja a un coche que va a 90 km/h.\\[3pt]
$t=\dfrac{30}{90-60}=\dfrac{30}{30}=\boxed{1\text{ h}}$
\end{exemple}
\end{minipage}

\begin{exemple}[Ejemplo: \ si no salen a la vez]
\begin{minipage}[c]{0.55\linewidth}
Un coche sale a las 9 h a 60 km/h. Otro sale a las 10 h del mismo lugar a 90 km/h. ¿Cuándo lo alcanza?
\end{minipage}\hfill
\begin{minipage}[c]{0.42\linewidth}
En 1 h, el primero saca \textbf{60 km}.\\[3pt]
$t=\dfrac{60}{90-60}=2$ h \ $\to$ \ \boxed{\text{a las 12 h}}
\end{minipage}
\end{exemple}

\begin{exemple}[Ejemplo: \ un ciclista va a 18 km/h. ¿Cuánto tarda en hacer 45 km?]
$t=\dfrac{d}{v}=\dfrac{45}{18}=2{,}5$ h $=\boxed{2\text{ h }30\text{ min}}$ \qquad \nota{¡0,5 h son 30 minutos, no 50!}
\end{exemple}

% ============================================================ 3.11
\clearpage
\seccio{sec11}{Repaso: proporcionalidad simple}{¿El doble, el doble? ¿O el doble, la mitad?}

\begin{tcbraster}[raster columns=2,raster equal height,raster column skip=10pt]
\begin{tcolorbox}[fmbase,colback=fmK!8,colframe=fmK,coltitle=white,title={\faArrowUp\ \ Directamente proporcionales},before upper=\centering]
El doble $\to$ el doble. \ Se conserva el \textbf{cociente}.\par\smallskip
\taularod[fmK]{>{\columncolor{fmK}\color{white}\heavy}l*{4}{w{c}{0.9cm}}}{
kg & 1 & 2 & 3 & 5\\
\rowcolor{fmK!8}\cellcolor{fmK}€ & 2,40 & 4,80 & 7,20 & 12
}\par\smallskip
$\dfrac{2{,}40}{1}=\dfrac{4{,}80}{2}=\dots=\cK{2{,}40}=k$
\end{tcolorbox}
\begin{tcolorbox}[fmbase,colback=fmB!8,colframe=fmB,coltitle=white,title={\faArrowDown\ \ Inversamente proporcionales},before upper=\centering]
El doble $\to$ la mitad. \ Se conserva el \textbf{producto}.\par\smallskip
\taularod[fmB]{>{\columncolor{fmB}\color{white}\heavy}l*{4}{w{c}{0.9cm}}}{
pintores & 1 & 2 & 3 & 4\\
\rowcolor{fmB!8}\cellcolor{fmB}días & 60 & 30 & 20 & 15
}\par\smallskip
$1\cdot60=2\cdot30=\dots=\cB{60}$
\end{tcolorbox}
\end{tcbraster}

\begin{center}
\taularod{>{\heavy}l<{\hspace{10pt}}>{\hspace{10pt}}c<{\hspace{10pt}}>{\hspace{10pt}}c<{\hspace{6pt}}}{
\rowcolor{fmblue} & \color{white}\heavy Directa (D) & \color{white}\heavy Inversa (I)\\
Regla de tres & $\left.\begin{array}{l}a\to b\\ c\to x\end{array}\right\}\ x=\dfrac{b\cdot c}{a}$ & $\left.\begin{array}{l}a\to b\\ c\to x\end{array}\right\}\ x=\dfrac{a\cdot b}{c}$\\[10pt]
\rowcolor{fmsurface}\rule[-15pt]{0pt}{36pt}Reducción a la unidad & $1\to\dfrac{b}{a}$, \ y multiplicamos por $c$ & $1\to a\cdot b$, \ y dividimos entre $c$
}
\end{center}

\begin{minipage}[t]{0.49\linewidth}
\begin{exemple}[Ejemplos]
3 kg cuestan 7,20 €. ¿Y 5 kg? \ \chip{fmK}{D}\\[2pt]
\hspace*{1em}$x=\dfrac{7{,}20\cdot5}{3}=\boxed{12\text{ €}}$\\[4pt]
4 pintores tardan 15 días. ¿Y 6 pintores? \ \chip{fmB}{I}\\[2pt]
\hspace*{1em}$x=\dfrac{4\cdot15}{6}=\boxed{10\text{ días}}$
\end{exemple}
\end{minipage}\hfill
\begin{minipage}[t]{0.48\linewidth}
\begin{compte}[¡No todas son proporcionales!]
La \textbf{edad} y la \textbf{altura} crecen a la vez, pero con el doble de edad no tenemos el doble de altura.\\[3pt]
El lado de un cuadrado y su área tampoco: el doble de lado, el área \textbf{cuádruple}.
\end{compte}
\end{minipage}

\begin{exemple}[Más ejemplos]
\begin{minipage}[t]{0.48\linewidth}
A 90 km/h tardamos 4 h. ¿Y a 120 km/h? \ \chip{fmB}{I}\\[3pt]
\hspace*{1em}$x=\dfrac{90\cdot4}{120}=\boxed{3\text{ h}}$
\end{minipage}\hfill
\begin{minipage}[t]{0.49\linewidth}
5 bolígrafos cuestan 3,50 €. ¿Y 8? \ \chip{fmK}{D} \ \nota{a la unidad}\\[3pt]
\hspace*{1em}1 boli: $3{,}50:5=0{,}70$ € \ $\to$ \ $8\cdot0{,}70=\boxed{5{,}60\text{ €}}$
\end{minipage}
\end{exemple}

% ============================================================ 3.12
\clearpage
\seccio{sec12}{Proporcionalidad compuesta}{Tres o más magnitudes: las comparamos de una en una.}

\begin{minipage}[c]{0.42\linewidth}\raggedright
\pas{1} Hacemos una \textbf{tabla} con los datos; la $x$ al final.\\[6pt]
\pas{2} Comparamos \textbf{cada} magnitud con la de la $x$, por separado: \chip{fmK}{D} o \chip{fmB}{I}?\\[6pt]
\pas{3} Multiplicamos el valor conocido por cada razón:\\[2pt]
\hspace*{1.6em}\chip{fmK}{D}: \ $\dfrac{\text{nuevo}}{\text{antiguo}}$\qquad \chip{fmB}{I}: \ $\dfrac{\text{antiguo}}{\text{nuevo}}$
\end{minipage}\hfill
\begin{minipage}[c]{0.55\linewidth}
\begin{exemple}[Ejemplo: \ 5 máquinas, 8 h al día, hacen 400 piezas. ¿Y 3 máquinas, 10 h al día?]
\centering
\taularod[fmex]{lccc}{
\rowcolor{fmex}\color{white}\heavy & \color{white}\heavy máquinas & \color{white}\heavy horas/día & \color{white}\heavy piezas\\
antes & 5 & 8 & 400\\
\rowcolor{fmexbg}ahora & 3 & 10 & $x$\\
 & \chip{fmK}{D} & \chip{fmK}{D} &
}\par\smallskip
$x=400\cdot\dfrac{3}{5}\cdot\dfrac{10}{8}=\boxed{300\text{ piezas}}$
\end{exemple}
\end{minipage}

\begin{exemple}[Ejemplo: \ 6 obreros, trabajando 8 h al día, tardan 10 días. ¿Cuántos días tardarán 4 obreros trabajando 6 h al día?]
\begin{minipage}[c]{0.5\linewidth}\centering
\taularod[fmex]{lccc}{
\rowcolor{fmex}\color{white}\heavy & \color{white}\heavy obreros & \color{white}\heavy horas/día & \color{white}\heavy días\\
antes & 6 & 8 & 10\\
\rowcolor{fmexbg}ahora & 4 & 6 & $x$\\
 & \chip{fmB}{I} & \chip{fmB}{I} &
}
\end{minipage}\hfill
\begin{minipage}[c]{0.47\linewidth}
\nota{menos obreros o menos horas $\to$ \textbf{más} días}\\[4pt]
$x=10\cdot\dfrac{6}{4}\cdot\dfrac{8}{6}=\boxed{20\text{ días}}$\\[2pt]
\nota{las razones inversas se invierten: antiguo entre nuevo}
\end{minipage}
\end{exemple}

\begin{exemple}[Ejemplo: \ 12 vacas comen 900 kg de pienso en 15 días. ¿Cuánto comen 20 vacas en 30 días?]
\begin{minipage}[c]{0.5\linewidth}\centering
\taularod[fmex]{lccc}{
\rowcolor{fmex}\color{white}\heavy & \color{white}\heavy vacas & \color{white}\heavy días & \color{white}\heavy kg\\
antes & 12 & 15 & 900\\
\rowcolor{fmexbg}ahora & 20 & 30 & $x$\\
 & \chip{fmK}{D} & \chip{fmK}{D} &
}
\end{minipage}\hfill
\begin{minipage}[c]{0.47\linewidth}
\nota{más vacas o más días $\to$ \textbf{más} pienso}\\[4pt]
$x=900\cdot\dfrac{20}{12}\cdot\dfrac{30}{15}=\boxed{3000\text{ kg}}$
\end{minipage}
\end{exemple}

\begin{minipage}[t]{0.5\linewidth}
\begin{truc}[Pregunta clave]
Para cada magnitud: \textbf{si esta aumenta} (y las demás no cambian), \textbf{¿la $x$ aumenta o disminuye?}
\end{truc}
\end{minipage}\hfill
\begin{minipage}[t]{0.47\linewidth}
\begin{idea}[También reduciendo a la unidad]
1 obrero, 1 hora al día: \ $10\cdot6\cdot8=480$ días.\\
4 obreros, 6 h: \ $480:(4\cdot6)=20$ días.
\end{idea}
\end{minipage}

% ============================================================ resum
\resumtitol
\begin{tcbraster}[raster columns=2,raster equal height=rows,raster column skip=14pt,raster row skip=22pt,raster before skip=14pt]
\rtile{3.1}{Porcentajes}{$p\pc=\dfrac{p}{100}$. \ $35\pc=0{,}35=\frac{7}{20}$.}
\rtile{3.2}{Cálculo}{$p\pc$ de $C=C\cdot\dfrac{p}{100}$. \ $10\pc$: entre 10.}
\rtile{3.3}{¿Qué \%?}{$p=\dfrac{P}{C}\cdot100$: la parte entre el total.}
\rtile{3.4}{Aumentos y descuentos}{$C_f=C_i\cdot I_V$, \ $I_V=1\pm\frac{p}{100}$.}
\rtile{3.5}{Cantidad inicial}{$C_i=\dfrac{C_f}{I_V}$. ¡No apliques el \% al final!}
\rtile{3.6}{Encadenadas}{Los $I_V$ se \textbf{multiplican}, no se suman.}
\rtile{3.7}{Interés compuesto}{$C_f=C_0\left(1+\frac{r}{100}\right)^t$}
\rtile{3.8}{Repartos}{Directo: $k=\dfrac{N}{a+b+c}$.\\[1ex] Inverso: a los inversos.}
\rtile{3.9}{Mezclas}{$p_m=\dfrac{m_1p_1+m_2p_2}{m_1+m_2}$}
\rtile{3.10}{Móviles}{$d=v\cdot t$. Se encuentran: $v_1+v_2$. Alcance: $v_1-v_2$.}
\rtile{3.11}{Proporcionalidad simple}{D: se conserva el cociente. I: el producto.}
\rtile{3.12}{Compuesta}{Cada magnitud por separado: D $\frac{\text{nuevo}}{\text{antiguo}}$,\\[1ex] I $\frac{\text{antiguo}}{\text{nuevo}}$.}
\end{tcbraster}
\end{document}
